\chapter{1977年湖南株洲卷（试点）}

\begin{problem}
\(y=2\sqrt{x^2}\) 和 \(y=2x\) 是不是同一个函数？为什么？
\end{problem}

\begin{solution}
两式的定义域都是 \(\R\)，但对任意 \(x\in\R\)，有
\[
2\sqrt{x^2}=2\lvert x\rvert=\begin{cases}
2x, & x\geq 0,\\
-2x, & x<0.
\end{cases}
\]
当 \(x=-1\) 时，\(2\sqrt{x^2}=2\)，而 \(2x=-2\)，两式的函数值不相等. 因此 \(y=2\sqrt{x^2}\) 与 \(y=2x\) 不是同一个函数.
\end{solution}

\begin{problem}
确定下列各式的符号：

\begin{enumerate}
\item \(1-\sin 800^\circ\)；
\item \(\cos 20^\circ-\sin 20^\circ\)；
\item \(\tan 288^\circ\cdot\tan 510^\circ\).
\end{enumerate}
\end{problem}

\begin{solution}
\begin{enumerate}
\item 因为 \(800^\circ=2\times 360^\circ+80^\circ\)，所以
\(1-\sin 800^\circ=1-\sin 80^\circ\). 又因为 \(0<\sin 80^\circ<1\)，故 \(1-\sin 800^\circ>0\).

\item \(\cos 20^\circ=\sin 70^\circ\)，所以
\(\cos 20^\circ-\sin 20^\circ=\sin 70^\circ-\sin 20^\circ\). 由于 \(0<20^\circ<70^\circ<90^\circ\)，正弦函数在该区间上递增，故 \(\sin 70^\circ-\sin 20^\circ>0\).

\item 由正切函数的周期性和诱导公式，
\[
\tan 288^\circ=\tan\left(360^\circ-72^\circ\right)=-\tan 72^\circ<0
\]
\[
\tan 510^\circ=\tan\left(360^\circ+150^\circ\right)=\tan 150^\circ=-\tan 30^\circ<0
\]
两个负数的乘积为正，所以 \(\tan 288^\circ\cdot\tan 510^\circ>0\).
\end{enumerate}
\end{solution}

\begin{problem}
已知 \(x\)，\(y\) 都是实数，且 \(\left(2x-1\right)^2+\left(y-8\right)^2\leq 0\)，求 \(\log_2xy\)？
\end{problem}

\begin{solution}
两个平方数都非负，故
\(\left(2x-1\right)^2+\left(y-8\right)^2\geq 0\). 结合题设不等式，只能有
\[
\left(2x-1\right)^2+\left(y-8\right)^2=0
\]
非负数之和为零时每一项都为零，因此
\(2x-1=0\)，\(y-8=0\)，即 \(x=\frac{1}{2}\)，\(y=8\). 于是 \(xy=4>0\)，对数有意义，且
\[
\log_2xy=\log_2\left(\frac{1}{2}\times 8\right)=\log_2 4=2
\]
\end{solution}

\begin{problem}
\(\sqrt{\left(\log_3 20\right)^2-6\log_3 20+9}\)？
\end{problem}

\begin{solution}
原式
\[
=\sqrt{\left(\log_3 20-3\right)^2}=\left\lvert\log_3 20-3\right\rvert
\]
因为底数 \(3>1\)，对数函数 \(\log_3x\) 递增，且
\(3=\log_3 27>\log_3 20\)，所以 \(\log_3 20-3<0\). 因此
\[
\left\lvert\log_3 20-3\right\rvert=3-\log_3 20
\]
\end{solution}

\begin{problem}
已知 \(\log_{1.5}7=a\)，\(\log_{1.5}8=b\)，求 \(\log_{56}1.5\)？
\end{problem}

\begin{solution}
因为 \(1.5>0\) 且 \(1.5\neq 1\)，可用换底公式. 又 \(56=7\times 8\)，所以
\[
\log_{56}1.5=\frac{\log_{1.5}1.5}{\log_{1.5}56}=\frac{1}{\log_{1.5}(7\times 8)}=\frac{1}{\log_{1.5}7+\log_{1.5}8}=\frac{1}{a+b}
\]
\end{solution}

\begin{problem}
化简：
\(\left(\frac{x^{-1}-1}{x^{-1}+1}-\frac{x^{-1}+1}{x^{-1}-1}\right)\cdot\left(2^{-1}-4^{-1}x^{-1}-4^{-1}x\right)\).
\end{problem}

\begin{solution}
原式有意义需满足 \(x\neq 0\)，\(x^{-1}+1\neq 0\)，\(x^{-1}-1\neq 0\)，即
\(x\in\R\mathbin{\backslash}\{-1,0,1\}\). 在此条件下，
\[
\frac{x^{-1}-1}{x^{-1}+1}-\frac{x^{-1}+1}{x^{-1}-1}
=\frac{1-x}{1+x}-\frac{1+x}{1-x}
=\frac{(1-x)^2-(1+x)^2}{1-x^2}
=\frac{-4x}{1-x^2}
\]
并且
\[
2^{-1}-4^{-1}x^{-1}-4^{-1}x
=\frac{1}{2}-\frac{1}{4x}-\frac{x}{4}
=\frac{2x-1-x^2}{4x}
=-\frac{(x-1)^2}{4x}
\]
所以原式
\[
=\frac{-4x}{1-x^2}\cdot\left(-\frac{(x-1)^2}{4x}\right)
=\frac{(x-1)^2}{1-x^2}
=\frac{(1-x)^2}{(1-x)(1+x)}
=\frac{1-x}{1+x}
\]
\end{solution}

\begin{problem}
\(h\) 为何数时，方程 \(2hx^2+\left(8h+1\right)x=-8h\) 有两个不相等的实根？
\end{problem}

\begin{solution}
将方程移项得
\[
2hx^2+\left(8h+1\right)x+8h=0
\]
要使它是关于 \(x\) 的二次方程，必须有 \(2h\neq 0\)，即 \(h\neq 0\). 此时判别式为
\[
\Delta=\left(8h+1\right)^2-4\times 2h\times 8h=16h+1
\]
方程有两个不相等的实根，当且仅当 \(\Delta>0\)，即 \(h>-\frac{1}{16}\). 再结合 \(h\neq 0\)，得到
\[
-\frac{1}{16}<h<0\quad\text{或}\quad h>0
\]
\end{solution}

\begin{problem}
证明恒等式：
\(\frac{2\sin\theta\cdot\sin\left(\frac{\pi}{2}-\theta\right)+\sin\left(\frac{3\pi}{2}+\theta\right)}{1+\sin\left(\pi+\theta\right)+\sin^2\left(\pi-\theta\right)-\sin^2\left(\frac{\pi}{2}+\theta\right)}=\tan\left(\frac{\pi}{2}-\theta\right)\).
\end{problem}

\begin{solution}
先考察左式的定义条件. 利用诱导公式，分母为
\[
1-\sin\theta+\sin^2\theta-\cos^2\theta
=2\sin^2\theta-\sin\theta
=\sin\theta\left(2\sin\theta-1\right)
\]
因此在左式有意义时，\(\sin\theta\neq 0\) 且 \(2\sin\theta-1\neq 0\). 在此条件下，左式的分子为
\[
2\sin\theta\cos\theta-\cos\theta
=\left(2\sin\theta-1\right)\cos\theta
\]
于是左式
\[
=\frac{\left(2\sin\theta-1\right)\cos\theta}{\sin\theta\left(2\sin\theta-1\right)}
=\frac{\cos\theta}{\sin\theta}
\]
另一方面，\(\sin\theta\neq 0\) 时
\[
\tan\left(\frac{\pi}{2}-\theta\right)=\frac{\cos\theta}{\sin\theta}
\]
所以左式等于右式，恒等式得证.
\end{solution}

\begin{problem}
\begin{enumerate}
\item 三个半径都是 \(R\) 的圆两两外切，求过三个切点的三个圆弧所围成图形的面积.
\item 已知菱形面积为 \(392\mathrm{~cm}^2\)，有一角为 \(150^\circ\)，求边长.
\end{enumerate}
\end{problem}

\begin{solution}
\begin{enumerate}
\item 设三个圆心为 \(A\)，\(B\)，\(C\)，三个切点分别为 \(D\)，\(E\)，\(F\). 由于三个圆半径相等且两两外切，\(AB=BC=CA=2R\)，所以 \(\triangle ABC\) 是边长为 \(2R\) 的等边三角形，且三个圆弧分别对应圆心角 \(60^\circ\) 的三个扇形.

由图可知，所求面积 \(s\) 等于 \(\triangle ABC\) 的面积减去三个扇形 \(ADF\)，\(BDE\)，\(CEF\) 的面积.

    \solutionfigure{\bitmapfigure[width=4.8cm]{img/1977/hunan_zhuzhou_pilot/q9_solution_fig1.png}}

因此
\[
s=\frac{1}{2}\cdot 2R\cdot 2R\cdot\sin 60^\circ-3\cdot\frac{\pi R^2}{6}
=\sqrt{3}R^2-\frac{\pi}{2}R^2
=\left(\sqrt{3}-\frac{\pi}{2}\right)R^2
\]

\item 菱形的邻角互补，所以若 \(\angle B=150^\circ\)，则 \(\angle A=30^\circ\). 以对角线分成的两个全等三角形计算面积，得到

    \solutionfigure{\bitmapfigure[width=5.6cm]{img/1977/hunan_zhuzhou_pilot/q9_solution_fig2.png}}

\[
392=2\left(\frac{1}{2}AB\cdot AD\cdot\sin A\right)
\]
因为菱形的四边相等，\(AB=AD\)，代入 \(\angle A=30^\circ\) 得
\[
392=AB^2\sin 30^\circ=\frac{1}{2}AB^2
\]
所以 \(AB^2=784\)，且边长 \(AB=\sqrt{784}=28\mathrm{~cm}\).

答：菱形边长为 \(28\mathrm{~cm}\).
\end{enumerate}
\end{solution}

\begin{problem}
在高 \(150\mathrm{~m}\) 的山顶上测得正南两船的俯角分别为 \(15^\circ\) 和 \(75^\circ\)，求两船的距离.
\end{problem}

\begin{solution}
设山顶为 \(P\)，山脚为 \(C\)，较远的船为 \(A\)，较近的船为 \(B\). 因为两船在正南方向，\(A\)，\(B\)，\(C\) 共线，且 \(PC=150\mathrm{~m}\). 俯角等于相应的仰角，所以
\(AC=150\cot 15^\circ\)，\(BC=150\cot 75^\circ=150\tan 15^\circ\).

    \solutionfigure{\bitmapfigure[width=5.8cm]{img/1977/hunan_zhuzhou_pilot/q10_solution_fig1.png}}

于是
\[
AB=AC-BC=150\left(\cot 15^\circ-\tan 15^\circ\right)
\]
利用 \(\cot\alpha=\frac{\cos\alpha}{\sin\alpha}\) 和 \(\tan\alpha=\frac{\sin\alpha}{\cos\alpha}\)，有
\[
\cot 15^\circ-\tan 15^\circ
=\frac{\cos^2 15^\circ-\sin^2 15^\circ}{\sin 15^\circ\cos 15^\circ}
=\frac{\cos 30^\circ}{\frac{1}{2}\sin 30^\circ}
=\frac{\frac{\sqrt{3}}{2}}{\frac{1}{4}}
=2\sqrt{3}
\]
因此
\[
AB=300\sqrt{3}\mathrm{~m}\approx 519.6\mathrm{~m}
\]

答：两船的距离为 \(300\sqrt{3}\mathrm{~m}\)，约为 \(519.6\mathrm{~m}\).
\end{solution}

\begin{problem}
有某种农药一桶，倒出 \(8\text{ 升}\) 后，用水补满；然后又倒出 \(4\text{ 升}\)，再用水补满. 于是测得桶中农药与水之比为 \(18:7\)，求木桶的容积.
\end{problem}

\begin{solution}
设木桶容积为 \(x\text{ 升}\). 由题意必须有 \(x>8\). 第一次倒出 \(8\text{ 升}\) 农药并补水后，桶内农药和水的体积分别为 \(x-8\) 和 \(8\)，所以混合液中农药、水的体积分数分别为 \(\frac{x-8}{x}\) 和 \(\frac{8}{x}\).

第二次倒出 \(4\text{ 升}\) 混合液，倒出的农药为 \(\frac{x-8}{x}\times 4\)，倒出的水为 \(\frac{8}{x}\times 4\). 再次补入 \(4\text{ 升}\) 水后，桶内农药和水的体积分别为
\[
x-8-\frac{x-8}{x}\times 4
=\frac{x^2-12x+32}{x}
\]
\[
8-\frac{8}{x}\times 4+4
=\frac{12x-32}{x}
\]
由农药与水之比为 \(18:7\)，得
\[
\frac{x^2-12x+32}{12x-32}=\frac{18}{7}
\]
交叉相乘并整理，得
\[
7x^2-300x+800=0
\]
因式分解得
\[
\left(7x-20\right)\left(x-40\right)=0
\]
所以 \(x=\frac{20}{7}\) 或 \(x=40\). 由于 \(x>8\)，舍去 \(\frac{20}{7}\)，取 \(x=40\).

答：桶的容积为 \(40\text{ 升}\).
\end{solution}

\begin{problem}
甲船在 \(A\) 点发现乙船在北偏东 \(60^\circ\) 的 \(B\) 处，乙船以每小时 \(a\text{ 里}\) 的速度向北行驶，已知甲船速度是每小时 \(\sqrt{3}a\text{ 里}\)，问甲船应依什么方向前进，才能最快地与乙船相会？
\end{problem}

\begin{solution}
设经过 \(t>0\) 小时两船在 \(C\) 点相遇. 乙船沿正北方向行驶，所以 \(BC=at\)；甲船的路程为 \(AC=\sqrt{3}at\). 由 \(AB\) 是北偏东 \(60^\circ\) 方向，且 \(BC\) 指向正北，得
\[
\angle B=180^\circ-60^\circ=120^\circ
\]

    \solutionfigure{\bitmapfigure[width=5.0cm]{img/1977/hunan_zhuzhou_pilot/q12_solution_fig1.png}}

令 \(L=AB>0\). 对任意可能的相遇时刻 \(t\)，乙船位于正北方向的点 \(C\)，且 \(BC=at\). 甲船在时刻 \(t\) 至多行驶 \(\sqrt{3}at\)，因此必须有 \(AC\leq\sqrt{3}at\). 由 \(\angle B=120^\circ\) 和余弦定理，
\[
AC^2=L^2+(at)^2-2L\cdot at\cos120^\circ=L^2+Lat+a^2t^2
\]
所以
\[
L^2+Lat+a^2t^2\leq 3a^2t^2
\]
整理得
\[
(2at+L)(at-L)\geq 0
\]
由于 \(a>0\)、\(t\geq 0\) 且 \(L>0\)，有 \(2at+L>0\)，故 \(at\geq L\)，从而任何相遇时间都满足 \(t\geq L/a\). 当 \(t=L/a\) 时取等号，甲船沿 \(AC\) 以全速直线行驶即可相遇，因而这是最短相遇时间. 此时 \(AB=BC=L\)，\(AC=\sqrt{3}L\). 在 \(\triangle ABC\) 中由正弦定理，
\[
\frac{AC}{\sin B}=\frac{BC}{\sin A}
\]
因此
\[
\sin A=\frac{BC\sin B}{AC}=\frac{L\sin120^\circ}{\sqrt{3}L}=\frac{1}{2}
\]
又因为 \(A+B+C=180^\circ\)，且 \(B=120^\circ\)，所以 \(0<A<60^\circ\)，从而 \(A=30^\circ\). 故甲船航线与正北方向的夹角为
\[
\angle NAC=60^\circ-\angle BAC=60^\circ-30^\circ=30^\circ
\]

答：甲船应以北偏东 \(30^\circ\) 的方向前进，才能最快地与乙船相会.
\end{solution}

\begin{problem}
双曲线中心在原点，焦点在坐标轴上；焦距为 \(14\)，顶点间距离为 \(12\)，求双曲线方程.
\end{problem}

\begin{solution}
双曲线的焦点在 \(x\) 轴或 \(y\) 轴上，需要分别讨论两种方向.

若焦点在 \(x\) 轴上，设方程为
\[
\frac{x^2}{a^2}-\frac{y^2}{b^2}=1
\]
顶点间距离为 \(2a=12\)，故 \(a=6\)；焦距为 \(2c=14\)，故 \(c=7\). 双曲线参数满足 \(c^2=a^2+b^2\)，所以
\[
b^2=c^2-a^2=7^2-6^2=13
\]
因此方程为
\[
\frac{x^2}{36}-\frac{y^2}{13}=1
\]

    \solutionfigure{\bitmapfigure[width=4.8cm]{img/1977/hunan_zhuzhou_pilot/q13_solution_fig1.png}}

若焦点在 \(y\) 轴上，设方程为
\[
\frac{y^2}{a^2}-\frac{x^2}{b^2}=1
\]
同样由 \(2a=12\)，\(2c=14\) 得 \(a=6\)，\(c=7\)，再由 \(c^2=a^2+b^2\) 得 \(b^2=13\). 因此方程为

    \solutionfigure{\bitmapfigure[width=3.2cm]{img/1977/hunan_zhuzhou_pilot/q13_solution_fig2.png}}

\[
\frac{y^2}{36}-\frac{x^2}{13}=1
\]
\end{solution}

\begin{problem}
从等腰 \(\triangle ABC\) 的顶点 \(A\) 引直线交 \(BC\) 于 \(D\)，交外接圆于 \(E\)，求证 \(AB\) 是 \(AD\) 和 \(AE\) 的比例中项.

\bitmapfigure[float=inline,width=5.0cm,trim=1.0cm 0 0 0,clip]{img/1977/hunan_zhuzhou_pilot/q14_fig1.png}
\end{problem}

\begin{solution}
连接 \(BE\). 因为 \(\triangle ABC\) 是等腰三角形，\(AB=AC\)，所以
\[
\angle ABC=\angle ACB
\]
又因为 \(A\)，\(B\)，\(C\)，\(E\) 共圆，\(\angle AEB\) 与 \(\angle ACB\) 都是所对弦 \(AB\) 所对的圆周角，故
\[
\angle AEB=\angle ACB
\]
由于 \(D\) 在 \(BC\) 上、\(E\) 在直线 \(AD\) 上，有
\[
\angle ABD=\angle ABC=\angle AEB
\]
并且 \(\angle BAD=\angle EAB\).
因此 \(\triangle ABD\sim\triangle AEB\). 对应边满足
\[
\frac{AB}{AE}=\frac{AD}{AB}
\]
从而 \(AB^2=AD\cdot AE\). 所以 \(AB\) 是 \(AD\) 和 \(AE\) 的比例中项.
\end{solution}

\begin{problem}
如果一个直角三角形三边的长成等差数列，试证明三边的比是 \(3:4:5\).
\end{problem}

\begin{solution}
设三边按从小到大的顺序为 \(a-d\)，\(a\)，\(a+d\)，其中公差 \(d>0\)，且 \(a-d>0\). 最大边是斜边，所以由勾股定理
\[
a^2+\left(a-d\right)^2=\left(a+d\right)^2
\]
移项并展开，得
\[
a^2=\left(a+d\right)^2-\left(a-d\right)^2
=a^2+2ad+d^2-a^2+2ad-d^2=4ad
\]
由于 \(a>0\)，可除以 \(4a\)，得到 \(d=\frac{a}{4}\). 因此三边之比为
\[
\left(a-\frac{a}{4}\right):a:\left(a+\frac{a}{4}\right)
=\frac{3a}{4}:a:\frac{5a}{4}=3:4:5
\]

    \solutionfigure{\bitmapfigure[width=4.4cm]{img/1977/hunan_zhuzhou_pilot/q15_solution_fig1.png}}
\end{solution}

\section{参考题}

\begin{problem}
试求 \(f(x)=x^4-42x^3+435x^2+113x-66\) 和 \(g(x)=x^2-x-462\) 的公根.
\end{problem}

\begin{solution}
先分解二次多项式：
\[
g(x)=x^2-x-462=(x-22)(x+21)
\]
所以公根只能是 \(22\) 或 \(-21\). 直接代入 \(f(x)\)，
\[
\begin{aligned}
f(22)&=22^4-42\times22^3+435\times22^2+113\times22-66\\
&=234256-447216+210540+2486-66=0,
\end{aligned}
\]
而
\[
\begin{aligned}
f(-21)&=(-21)^4-42\times(-21)^3+435\times(-21)^2+113\times(-21)-66\\
&=194481+388962+191835-2373-66=772839\neq 0.
\end{aligned}
\]
因此 \(22\) 是两个多项式的公根，\(-21\) 不是，所求公根为 \(x=22\).
\end{solution}

\begin{problem}
求函数 \(y=\left(2x+1\right)^{\frac{1}{2}}\left(-\frac{1}{3}x+1\right)^3\) 的导数.
\end{problem}

\begin{solution}
函数的定义域满足 \(2x+1\geq 0\)，即 \(x\geq-\frac{1}{2}\). 在可导区间 \(x>-\frac{1}{2}\) 上，令
\[
u=\left(2x+1\right)^{\frac{1}{2}}
\]
令
\[
v=\left(1-\frac{x}{3}\right)^3
\]
由链式法则，
\[
u'=\frac{1}{2}\left(2x+1\right)^{-\frac{1}{2}}\cdot 2=\left(2x+1\right)^{-\frac{1}{2}}
\]
并且
\[
v'=3\left(1-\frac{x}{3}\right)^2\cdot\left(-\frac{1}{3}\right)=-\left(1-\frac{x}{3}\right)^2
\]
再由乘积求导法则，
\[
y'=u'v+uv'
=\left(1-\frac{x}{3}\right)^3\left(2x+1\right)^{-\frac{1}{2}}-\left(2x+1\right)^{\frac{1}{2}}\left(1-\frac{x}{3}\right)^2
\]
\end{solution}

\begin{problem}
有一窗户，顶部为一抛物线的拱形，窗宽 \(120\mathrm{~cm}\)，边高 \(150\mathrm{~cm}\)，顶点高 \(180\mathrm{~cm}\)，求它的面积.
\end{problem}

\begin{solution}
取顶点为原点，水平方向为 \(x\) 轴，竖直向上为 \(y\) 轴. 设抛物线方程为
\[
x^2=-2py
\]
窗宽为 \(120\mathrm{~cm}\)，故右侧边的上端点坐标为 \(A(60,-30)\)，因为顶点比边高 \(180-150=30\mathrm{~cm}\). 代入抛物线方程，得
\[
60^2=-2p\times(-30)
\]
所以 \(p=60\).
所以抛物线方程为
\[
x^2=-120y
\]
即 \(y=-\frac{x^2}{120}\).

    \solutionfigure{\bitmapfigure[width=5.8cm,trim=0 0 0 1.2cm,clip]{img/1977/hunan_zhuzhou_pilot/q18_solution_fig1.png}}

图中右半边阴影面积是曲线与 \(x\) 轴之间的面积，故
\[
\int_0^{60}\frac{x^2}{120}\,\mathrm{d}x
=\left.\frac{x^3}{360}\right|_0^{60}
=600\mathrm{~cm}^2
\]
整个窗户可看作宽为 \(120\mathrm{~cm}\)、高为 \(180\mathrm{~cm}\) 的矩形减去左右两个全等阴影区域，因此
\[
S=2\left(180\times 60-600\right)\mathrm{~cm}^2
=20400\mathrm{~cm}^2
=2.04\mathrm{~m}^2
\]

答：所求窗户的面积是 \(2.04\mathrm{~m}^2\).
\end{solution}
